\subsection{The resolved tail repair with its prescribed edge boundary}

\begin{theorem}[AT17: Proved; the relative resolved-tail construction]
\label{AT17}
Keep the coefficient ring, targets, coefficient module
$L=I_N\langle s\rangle$, cutoff $W=N+s$, integer $m$,
and finite coefficient resolution $W^\bullet$ of AT16.
Let $X\to A,B$ be the stated actual maps, and assume
$X_0=S_R$. Suppose the simplex-one coefficient map
has a specified weight-preserving $S_R$-linear cochain
contraction
\[
 dH+Hd=1-ir,\qquad
 i:S_R\to X_1,\quad r:X_1\to S_R,\quad ri=1.
 \tag{AT.83}
\]
This data is supplied by the even-to-odd disk
decomposition of the refined C21 simplex-one algebra.

Now take actual $S_R$-linear simplicial data on
the disk parameter, rather than the sphere:
\[
 z:L\otimes_R D(2)\to Z^4(X),\qquad
 \eta:L\otimes_R D(2)\to X^3,\qquad d\eta=mz.
 \tag{AT.84}
\]
There is an explicitly defined relative resolved
strict PD extension $X\hookrightarrow Y_{\mathrm{rel}}^{\mathrm{tf}}$
with actual extending maps to $A,B$.
Its terms are flat, all weights below $W$ are
unchanged, and its entire simplex-zero and
simplex-one algebras are $X_0$ and $X_1$.
Its rationalized inclusion is a quasi-isomorphism
in every simplex and weight, and its coefficient
$H^0$ is $S_R$ in every simplex and weight.
The top $Q_0$ generators supply vertical primitives
for all values of the prescribed disk cycle $z$.

No assertion about its nonlinear positive integral
cohomology is made.
\end{theorem}

\begin{proof}
Write $z_1,\eta_1$ for the restrictions of
\textup{(AT.84)} to $L\otimes K(1)$.
In simplex one put
\[
                    a=H z_1,\qquad
                    c=H(ma-\eta_1).
 \tag{AT.85}
\]
The first argument is closed in form degree four.
Equation \textup{(AT.83)} therefore gives $da=z_1$.
The second argument is closed in degree three,
since $d(ma-\eta_1)=mz_1-mz_1=0$.
Thus $dc=ma-\eta_1$.
The outputs have positive form degrees three and two.
Their two faces in $X_0=S_R$ are therefore zero.
Consequently they specify actual maps on the
simplicial $K(1)$ parameter, by Dold--Kan.
They are $S_R$-linear and preserve every weight
because $H$ does. No $S_R$-linear horizontal
primitive on a coefficient tail has been chosen.

Use the complex $P$ and its truncation $W^\bullet$
from \textup{(AT.75)--(AT.76)}.
Use the twist \textup{(AT.77)} with the present
disk-parameter $z,\eta$.
On the edge subparameter define generator values
\[
 \phi(u_0)=a\varepsilon,\qquad
 \phi(v_0)=c\varepsilon,\qquad
 \phi(u_i)=\phi(v_i)=0\ (i>0),\qquad
 \phi(W^0)=0.
 \tag{AT.86}
\]
They satisfy $d_X\phi=\phi d_W+\delta|_{K(1)}$.
For $u_0,v_0$ these are the two equations just
proved for $a,c$. On all remaining terms they
are $\varepsilon\partial=0$.
Thus they give actual compatible boundary values
for the entire resolved generator complex.

Take the free relative PD algebra over $X$ on
$W^\bullet\otimes D(2)$ with differential
$d_W+\delta$, and impose the PD relations
\[
                 t|_{K(1)}=\phi(t)
                   \qquad(t\in W^\bullet).
 \tag{AT.87}
\]
The preceding cochain identity proves differential
stability of this relation ideal; all its other
compatibilities follow from the specified homogeneous
simplicial linear maps.
This defines $Y_{\mathrm{rel}}$.
In each simplex the inclusion $K(1)\subset D(2)$
has a free $R$-linear complement.
The free PD universal property therefore identifies
the underlying graded algebra in that simplex with
a free relative PD extension of $X$ on
$W^\bullet\otimes K(2)$.
This supplies a graded retraction in each simplex,
proves injectivity of $X$, and proves the lower-weight
and simplex-zero/one assertions.
A simplicial choice of these complements is neither
needed nor asserted.

We construct the actual target maps with their
prescribed edge values. Work first on the full
free complex $P$.
For each of its free coefficient generators in
form degrees zero through three, the prescribed
edge value $f\phi$ is a normalized edge cycle
in the corresponding target row.
Horizontal acyclicity supplies a normalized
two-simplex filling of that value.
Choose these fillings independently on the free
generators. They give a weight-preserving graded
simplicial map
\[
 h^{\mathrm{pre}}:P\otimes D(2)\to A
\]
of degree zero, whose edge restriction is
$f\phi$ pulled back from $W$.
This preliminary map is not asserted to commute
with differentials. All required rows are within
the degree-four bounds of AT16.

Its differential error
\[
                 E=f\delta-Dh^{\mathrm{pre}}
 \tag{AT.88}
\]
is a degree-one Hom cocycle.
It restricts to zero on $P\otimes K(1)$, because
the prescribed boundary values satisfy their
cochain identity. Thus it is an actual map from
$P\otimes K(2)$.
The target calculation \textup{(AT.79)} gives
$H^1\operatorname{Hom}(P\otimes K(2),A)=0$,
with its displayed interpretation as
$\operatorname{Ext}^2(L/mL,S_R)=0$.
Choose a degree-zero correcting cochain $k$
on this quotient with $Dk=E$.
Then $h=h^{\mathrm{pre}}+k$ obeys
$Dh=f\delta$ and retains exactly the prescribed
edge values.

The terminal equation on $P^{-1}$ is again
$h_0d_P=0$. It makes $h_0$ factor through the
actual quotient $W^0$.
We have therefore constructed generator values
on $W^\bullet\otimes D(2)$ satisfying both the
differential and every relation in
\textup{(AT.87)}.
The free PD universal property gives
$Y_{\mathrm{rel}}\to A$.
Repeat the same calculation for $B$.
The coefficient resolution was free at the only
stage where horizontal primitives were chosen;
no lifting through an unspecified nonfree-tail
Hom functor has entered this argument.

For rational exactness, first use the explicit
contraction of $W_{\mathbb Q}$ in \textup{(AT.80)}
and the translation $t'=t-a'(t)$, where
$a'=-\delta H$, exactly as in AT16.
This untwists the differential.
Its boundary value is now the cochain map
\[
               \psi=\phi-a'|_{K(1)}
              :W_{\mathbb Q}\otimes K(1)\to X_{\mathbb Q}.
\]
Fix one simplex.
Choose a splitting of $K(1)_n\subset D(2)_n$
as free $R$-modules and extend $\psi_n$ by zero
on that complement.
Because the splitting acts on the parameter and
has no vertical differential, this extension
is an actual $S_{\mathbb Q}$-linear cochain map.
Translate the rational generators once more by
that extension. Their boundary values become zero,
and their differential remains $d_W$.
In this simplex we have thus exhibited an actual
PD-DGA isomorphism with
\[
 X_{n,\mathbb Q}\otimes_{S_{\mathbb Q}}
    \operatorname{FreePD}_{S_{\mathbb Q}}
                  (W_{\mathbb Q}\otimes K(2)_n).
 \tag{AT.89}
\]
The rational arity contraction of AT16 applies.
It proves that the original inclusion of $X$
is a rational quasi-isomorphism in every simplex
and weight. The separate splittings are used
only to verify that actual inclusion; they are
not presented as a simplicial splitting.

Finally take the image of $Y_{\mathrm{rel}}$
in its rationalization.
Its torsion ideal is differential and PD-stable,
and both actual maps to torsionfree targets
factor through it.
The simplexwise graded retractions above show
that the torsionfree $X$ still embeds.
Flatness over the PID $R$ follows.
All lower weights and the first two simplex
algebras stay unchanged.
Rational coefficient exactness in degree zero,
together with the actual target coefficient
comparison above the cutoff, proves $H^0=S_R$
by the same coefficient-integrality argument
as in AT16.
This proves every assertion.
\end{proof}

There is also the actual augmentation comparison
with the relative two-term tail algebra.
Use copies of $L\otimes D(2)$ in forms two and
three, with $du=z$, $dv=mu-\eta$, and impose
$u|_{K(1)}=a$, $v|_{K(1)}=c$.
The coefficient-resolution augmentation on $u_0,v_0$,
zero on all higher generators and the bottom kernel,
gives a surjective strict PD-DGA map from
$Y_{\mathrm{rel}}^{\mathrm{tf}}$ to its torsionfree
quotient. The equations \textup{(AT.86)} show that
this map respects the boundary relations as well
as the differential.
Its rational quasi-isomorphism follows from the
same actual simplexwise translations.
The construction supplies this morphism, not an
integral low-cohomology equivalence of its source
and target.

